% Nguồn SGK Toán 9 Tập một — Bài 7, trang 44--48:
% https://cdn3.olm.vn/upload/thuvienso/4491668113-page-44-1773022935181.png
% https://cdn3.olm.vn/upload/thuvienso/4491668113-page-45-1773022935481.png
% https://cdn3.olm.vn/upload/thuvienso/4491668113-page-46-1773022935831.png
% https://cdn3.olm.vn/upload/thuvienso/4491668113-page-47-1773022936159.png
% https://cdn3.olm.vn/upload/thuvienso/4491668113-page-48-1773022936453.png
% Nguồn SBT Toán 9 Tập một — Bài 7, PDF trang 30--32:
% SBT TOAN 9 KNTT TAP 1.pdf
% SHA-256: d7ffd04618456682795977e5c47de7874161a35e8606e5bac1906881603a72f7

\section{Căn bậc hai và căn thức bậc hai}

\begin{lesson}
    \subsection{Kiến thức trọng tâm}

    \textbf{Mở đầu}

    Trong Vật lí, quãng đường $S$ (tính bằng mét) của một vật rơi tự do được cho bởi công thức $S=4{,}9t^2$, trong đó $t$ là thời gian rơi (tính bằng giây). Hỏi sau bao nhiêu giây thì vật sẽ chạm đất nếu được thả rơi tự do từ độ cao $122{,}5$ mét?

    \answerline

    \subsubsection{Căn bậc hai}

    \textbf{Tìm hiểu khái niệm căn bậc hai}

    \textbf{HĐ1.} Tìm các số thực $x$ sao cho $x^2=49$.

    \answerline

    \begin{keyconcept}
        Căn bậc hai của số thực không âm $a$ là số thực $x$ sao cho $x^2=a$.
    \end{keyconcept}

    \textbf{Nhận xét}
    \begin{itemize}
        \item Số âm không có căn bậc hai;
        \item Số $0$ có một căn bậc hai duy nhất là $0$;
        \item Số dương $a$ có đúng hai căn bậc hai đối nhau là $\sqrt{a}$ (căn bậc hai số học của $a$) và $-\sqrt{a}$.
    \end{itemize}

    \textbf{Ví dụ 1.} Tìm căn bậc hai của $81$.

    \textbf{Giải}

    Ta có $\sqrt{81}=9$ nên $81$ có hai căn bậc hai là $9$ và $-9$.

    \textbf{Luyện tập 1.} Tìm căn bậc hai của $121$.

    \answerline

    \textbf{Tính căn bậc hai của một số bằng máy tính cầm tay}

    Để tính các căn bậc hai của một số $a>0$, chỉ cần tính $\sqrt{a}$. Có thể dễ dàng làm điều này bằng cách sử dụng MTCT.

    \textbf{Ví dụ 2.} Sử dụng MTCT, tính căn bậc hai của $11{,}1$ (làm tròn đến chữ số thập phân thứ hai).

    \textbf{Giải}

    Bấm các phím thích hợp, màn hình hiện kết quả là $3{,}33166625$. Làm tròn kết quả đến chữ số thập phân thứ hai ta được $\sqrt{11{,}1}\approx3{,}33$.

    Vậy căn bậc hai của $11{,}1$ (làm tròn đến chữ số thập phân thứ hai) là $3{,}33$ và $-3{,}33$.

    \textbf{Luyện tập 2.} Sử dụng MTCT tìm căn bậc hai của $\dfrac{7}{11}$ (làm tròn đến chữ số thập phân thứ hai).

    \answerline

    \textbf{Tính chất của căn bậc hai}

    \textbf{HĐ2.} Tính và so sánh $\sqrt{a^2}$ và $|a|$ trong mỗi trường hợp sau:
    \begin{enumerate}[label=\alph*)]
        \item $a=3$;
        \item $a=-3$.
    \end{enumerate}

    \answerline

    \begin{keyconcept}
        \[
            \sqrt{a^2}=|a|
        \]
        với mọi số thực $a$.
    \end{keyconcept}

    \textbf{Ví dụ 3.} Không sử dụng MTCT, tính:
    \begin{enumerate}[label=\alph*)]
        \item $1-\sqrt{2}+\sqrt{(1+\sqrt{2})^2}$;
        \item $\sqrt{(-3)^2}+3$.
    \end{enumerate}

    \textbf{Giải}

    a) Ta có $\sqrt{(1+\sqrt{2})^2}=|1+\sqrt{2}|=1+\sqrt{2}$ nên
    \[
        1-\sqrt{2}+\sqrt{(1+\sqrt{2})^2}
        =1-\sqrt{2}+(1+\sqrt{2})=2.
    \]

    b) Ta có $\sqrt{(-3)^2}=|-3|=3$ nên
    \[
        \sqrt{(-3)^2}+3=3+3=6.
    \]

    \textbf{Luyện tập 3.}
    \begin{enumerate}[label=\alph*)]
        \item Không sử dụng MTCT, tính: $\sqrt{6^2}$; $\sqrt{(-5)^2}$; $\sqrt{5}-\sqrt{(\sqrt{5}-1)^2}$.

        \answerline
        \item So sánh $3$ với $\sqrt{10}$ bằng hai cách:
        \begin{itemize}
            \item Sử dụng MTCT;
            \item Sử dụng tính chất của căn bậc hai số học đã học ở lớp 7: Nếu $0\le a<b$ thì $\sqrt{a}<\sqrt{b}$.
        \end{itemize}

        \answerline
    \end{enumerate}

    \subsubsection{Căn thức bậc hai}

    \textbf{Căn thức bậc hai}

    \textbf{HĐ3.} Viết biểu thức tính độ dài cạnh huyền $BC$ của tam giác vuông $ABC$, biết $AB=3$ cm và $AC=x$ cm.

    % TODO(HINH-NGUON): bo sung asset/crop dung tam giac ABC tu SGK trang 46; khong redraw xap xi khi chua doi chieu duoc vi tri/marker cua hinh goc.
    \answerline

    Biểu thức nhận được trong HĐ3 có dạng $\sqrt{A}$, trong đó $A$ là một biểu thức đại số. Tổng quát, ta có định nghĩa:

    \begin{keyconcept}
        Căn thức bậc hai là biểu thức có dạng $\sqrt{A}$, trong đó $A$ là một biểu thức đại số. $A$ được gọi là biểu thức lấy căn hoặc biểu thức dưới dấu căn.
    \end{keyconcept}

    \textbf{HĐ4.} Cho biểu thức $C=\sqrt{2x-1}$.
    \begin{enumerate}[label=\alph*)]
        \item Tính giá trị của biểu thức tại $x=5$.
        \item Tại $x=0$ có tính được giá trị của biểu thức không? Vì sao?
    \end{enumerate}

    \answerline
    \answerline

    \begin{keyconcept}
        $\sqrt{A}$ xác định khi $A$ lấy giá trị không âm và ta thường viết là $A\ge0$. Ta nói $A\ge0$ là điều kiện xác định (hay điều kiện có nghĩa) của $\sqrt{A}$.
    \end{keyconcept}

    \textbf{Ví dụ 4.} Xét căn thức $\sqrt{2x+1}$.
    \begin{enumerate}[label=\alph*)]
        \item Tìm điều kiện xác định của căn thức.
        \item Tính giá trị của căn thức đã cho tại $x=0$ và $x=4$.
    \end{enumerate}

    \textbf{Giải}

    a) Điều kiện xác định của căn thức là $2x+1\ge0$ hay $x\ge-\dfrac{1}{2}$.

    b) Tại $x=0$ (thỏa mãn điều kiện xác định) căn thức có giá trị là $\sqrt{2\cdot0+1}=1$.

    Tại $x=4$ (thoả mãn điều kiện xác định) căn thức có giá trị là $\sqrt{2\cdot4+1}=3$.

    \textbf{Luyện tập 4.} Cho căn thức $\sqrt{5-2x}$.
    \begin{enumerate}[label=\alph*)]
        \item Tìm điều kiện xác định của căn thức.
        \item Tính giá trị của căn thức tại $x=2$.
    \end{enumerate}

    \answerline
    \answerline

    \textbf{Hằng đẳng thức $\sqrt{A^2}=|A|$}

    Tương tự như căn bậc hai của một số thực không âm, với $A$ là một biểu thức, ta cũng có:
    \begin{itemize}
        \item Với $A\ge0$ ta có $\sqrt{A}\ge0$; $(\sqrt{A})^2=A$;
        \item $\sqrt{A^2}=|A|$.
    \end{itemize}

    \textbf{Ví dụ 5.} Rút gọn các biểu thức sau:
    \begin{enumerate}[label=\alph*)]
        \item $\sqrt{(1-x)^2}$ với $x<0$;
        \item $\sqrt{1-2x+x^2}$ với $x>2$.
    \end{enumerate}

    \textbf{Giải}

    a) Từ giả thiết $x<0$ suy ra $1-x>0$. Do đó
    \[
        \sqrt{(1-x)^2}=1-x.
    \]

    b) Áp dụng hằng đẳng thức bình phương của một hiệu và hằng đẳng thức $\sqrt{A^2}=|A|$, ta có
    \[
        \sqrt{1-2x+x^2}=\sqrt{(1-x)^2}=|1-x|.
    \]
    Do giả thiết $x>2$ suy ra $1-x<0$ nên $|1-x|=-(1-x)=x-1$. Vì vậy
    \[
        \sqrt{1-2x+x^2}=\sqrt{(1-x)^2}=x-1
    \]
    với $x>2$.

    \textbf{Luyện tập 5.}
    \begin{enumerate}[label=\alph*)]
        \item Rút gọn biểu thức $x\sqrt{x^6}$ ($x<0$).

        \answerline
        \item Rút gọn và tính giá trị của biểu thức $x+\sqrt{4x^2-4x+1}$ tại $x=-2{,}5$.

        \answerline
        \answerline
    \end{enumerate}

    \textbf{Vận dụng.} Trở lại tình huống mở đầu.
    \begin{enumerate}[label=\alph*)]
        \item Viết công thức tính thời gian $t$ (giây) cần thiết để vật rơi được quãng đường $S$ (mét).
        \item Sử dụng công thức tìm được trong câu a, hãy trả lời câu hỏi trong tình huống mở đầu.
    \end{enumerate}

    \answerline
    \answerline
\end{lesson}

\begin{exercises}
    \subsection{Bài tập}

    \begin{questions}[series=SGK]
        \item Tìm căn bậc hai của mỗi số sau (làm tròn đến chữ số thập phân thứ hai):
        \begin{questions}
            \item $24{,}5$;
            \item $\dfrac{9}{10}$.
        \end{questions}

        \answerline

        \item Để chuẩn bị trồng cây trên vỉa hè, người ta để lại những ô đất hình tròn có diện tích khoảng $2\,\text{m}^2$. Em hãy ước lượng (với độ chính xác $0{,}005$) đường kính của các ô đất đó khoảng bao nhiêu mét?

        \answerline
        \answerline

        \item Tìm điều kiện xác định của $\sqrt{x+10}$ và tính giá trị của căn thức tại $x=-1$.

        \answerline
        \answerline

        \item Tính:
        \[
            \sqrt{5{,}1^2};\qquad
            \sqrt{(-4{,}9)^2};\qquad
            -\sqrt{(-0{,}001)^2}.
        \]

        \answerline

        \item Rút gọn các biểu thức sau:
        \begin{questions}
            \item $\sqrt{(2-\sqrt{5})^2}$;
            \item $3\sqrt{x^2}-x+1$ ($x<0$);
            \item $\sqrt{x^2-4x+4}$ ($x<2$).
        \end{questions}

        \answerline
        \answerline

        \item Không dùng MTCT, chứng tỏ biểu thức $A$ có giá trị là số nguyên:
        \[
            A=\sqrt{(1+2\sqrt{2})^2}-\sqrt{(1-2\sqrt{2})^2}.
        \]

        \answerline
        \answerline
    \end{questions}
\end{exercises}

\begin{practice}
    \subsection{Luyện tập}

    \begin{questions}[series=SBT]
        \item Tìm các căn bậc hai của mỗi số sau:
        \begin{questions}
            \item $81$;
            \item $\dfrac{16}{625}$;
            \item $0{,}0121$;
            \item $6\,400$.
        \end{questions}

        \answerline

        \item Sử dụng MTCT tính:
        \begin{questions}
            \item $\sqrt{17}$ (làm tròn kết quả đến chữ số thập phân thứ ba);
            \item Các căn bậc hai của $4\,021$ (làm tròn kết quả đến hàng phần trăm);
            \item Giá trị biểu thức $\dfrac{-11+\sqrt{11^2-4\cdot3\cdot2}}{2\cdot3}$ (làm tròn kết quả với độ chính xác $0{,}005$).
        \end{questions}

        \answerline
        \answerline

        \item Rút gọn:
        \begin{questions}
            \item $(\sqrt{4{,}1})^2-(-\sqrt{6{,}1})^2$;
            \item $(\sqrt{101})^2-\sqrt{(-99)^2}$;
            \item $\sqrt{(\sqrt{3}+2\sqrt{2})^2}-(-\sqrt{3}+2\sqrt{2})$;
            \item $\sqrt{(\sqrt{10}+3)^2}-\sqrt{(\sqrt{10}-3)^2}$.
        \end{questions}

        \answerline
        \answerline

        \item Sử dụng các hằng đẳng thức đáng nhớ hiệu hai bình phương và bình phương của một hiệu, rút gọn:
        \begin{questions}
            \item $(\sqrt{3}+\sqrt{2})(\sqrt{3}-\sqrt{2})$;
            \item $\sqrt{2-2\sqrt{2}+1}$.
        \end{questions}

        \answerline
        \answerline

        \item Khi giải phương trình $ax^2+bx+c=0$ ($a$, $b$, $c$ là ba số thực đã cho, $a\ne0$), ta phải tính giá trị của căn thức bậc hai $\sqrt{b^2-4ac}$. Hãy tính giá trị của căn thức này với các phương trình sau:
        \begin{questions}
            \item $x^2+5x+6=0$;
            \item $4x^2-5x-6=0$;
            \item $-3x^2-2x+33=0$.
        \end{questions}

        \answerline
        \answerline

        \item Rút gọn các biểu thức sau:
        \begin{questions}
            \item $\sqrt{49x^4}-3x^2$;
            \item $\sqrt{a^6(a-b)^2}:(a-b)$ với $a<b<0$.
        \end{questions}

        \answerline
        \answerline

        \item Rút gọn rồi tính giá trị biểu thức $\sqrt{25(4x^2-4x+1)^2}$ tại $x=\sqrt{3}$.

        \answerline
        \answerline
    \end{questions}
\end{practice}
